\documentclass[10pt,a4paper]{amsart}
\usepackage[margin=19mm]{geometry}
\usepackage{amsmath,amssymb,mathtools}
\usepackage[scr=rsfs,cal=euler]{mathalfa}
\usepackage{xfrac}
\usepackage[unicode,hidelinks,bookmarks=false]{hyperref}
\newcommand{\ie}{i.e.\ }
\newcommand{\cf}{cf.\ }
\newcommand{\resp}{respectively\ }
\newcommand{\Subsection}{\subsection}
\providecommand{\nobreakdash}{\nobreak\hbox{-}}
\setlength{\emergencystretch}{2em}
\multlinegap0pt
\usepackage{bm}
\usepackage{bbm}
\usepackage{dsfont}
\usepackage{xspace}
\usepackage{cancel}
\usepackage{mathtools}
\usepackage{amsthm}
\makeatletter
\@ifundefined{thm}{\newtheorem{thm}{Thm}}{}
\@ifundefined{prop}{\newtheorem{prop}[thm]{Proposition}}{}
\makeatother
\makeatletter
\expandafter\ifx\csname cases\endcsname\relax
\renewenvironment{cases}[1][l]{\matrix@check\cases\env@cases{#1}}{\endarray\right.}
\fi
\makeatother
\title[Transition asymptotics for the real solutions of the sinh-Go]{Transition asymptotics for the real solutions of the sinh-Gordon Painlev\'e III equation}
\author{Kenta Miyahara, Maxim L. Yattselev}
\date{Source: arXiv:2606.28579v1 (CC BY 4.0)}
\begin{document}
\maketitle

\noindent\emph{Source and licence.} Transition asymptotics for the real solutions of the sinh-Gordon Painlev\'e III equation. Version arXiv:2606.28579v1 (CC BY 4.0), distributed under the Creative Commons Attribution 4.0 International licence, as recorded in the arXiv licence metadata for this exact version and shown on the abstract page https://arxiv.org/abs/2606.28579v1. Full rights notice: licenses/arxiv-2606.28579-CC-BY-4.0.txt.\par\smallskip

\noindent\textit{Editorial scope.} Counted unit: the Prop whose source label is \texttt{prop:det} in the article (source0.tex lines 4568--4578, source label \texttt{prop:det}), delivered together with the article's own complete original proof of it (source0.tex lines 4579--4641, one \texttt{proof} environment of 63 lines). No mathematics is written, repaired or re-derived: statement and proof are the article's own text line for line. Required context retained uncounted: bimonial identity (lines 4561--4565). Every definition, display and label of those lines is retained unchanged. Omitted from this package and not redistributed: 1--4560, 4566--4567, 4642--4914. Nothing inside a retained excerpt is omitted or altered: every retained body is delivered exactly as the article's lines are written, so no omission receipt of any kind accompanies this package. Citations: the retained excerpts carry no citation command at all, so no bibliography entry of the article is transcribed and no reference is redistributed. Reference adaptation: none was needed and none was made; every reference inside the delivered unit resolves inside it, so no unresolved reference or citation is produced. Printed slips kept literal: no printed word, symbol, hypothesis or number of the counted statement or of its proof was altered. Layout and class: the article's own class is \texttt{amsart} with packages including amsmath, amsthm, amssymb, stmaryrd, url, here, xy, comment, mathtools, ascmac, ulem, framed, graphicx, nicematrix; the wrapper is the lane's pinned \texttt{amsart} editorial wrapper at 10pt on A4, which restates the theorem-like environments the retained text needs and adds the article's own macro definitions that the retained text needs (none), with extra wrapper packages (bm, bbm, dsfont, xspace, cancel, mathtools). The delivered numbering is the unit's own. Assets: no retained span contains a figure environment, a graphics-inclusion command, a PDF-page inclusion, a table or a TikZ picture; no image file is copied into this package, the package declares no asset dependency and no image file is redistributed with it. Licence: version arXiv:2606.28579v1 of this article is distributed under the Creative Commons Attribution 4.0 International licence, as recorded in the arXiv licence metadata for this exact version and shown on the abstract page https://arxiv.org/abs/2606.28579v1. A full rights notice is delivered at licenses/arxiv-2606.28579-CC-BY-4.0.txt. Characters: every retained excerpt is pure ASCII, so the delivered engine, XeLaTeX with its default Latin Modern text font, reports no missing character.
\begin{align}
    \label{bimonial identity}
    \sum_{k=0}^m\binom{m+n}k x^k y^{m-k} & = \sum_{k=0}^m \binom{-n}{k} (-x)^k(x+y)^{m-k} \\
    & = \sum_{k=0}^m (-1)^{n-1}\binom{-k-1}{n-1} x^k(x+y)^{m-k},
\end{align}

\begin{prop}
\label{prop:det}
Let \( T=[t_{j-i}]_{i,j=1}^N \) and \( H=[h_{i+j-1}]_{i,j=1}^N \) be a Toeplitz and a Hankel matrix, respectively, with \( t_k=0 \) when \( k<0 \) and \( h_k=0 \) when \( k>N \) (in all matrices, \( i \) stands for the row index and \( j \) for the column one). Further, let
\[
L = \left[\left(\frac\ell2\right)^{i+j-1}(-1)^{i-1}\binom{-j}{i-1}\right]_{i,j=1}^N
\]
for some constant \( \ell \). Assume that \( \det(T-HL)\neq 0 \). If \( \vec v=(v_1,\ldots,v_N)^\mathsf{T} \) is the solution of \( (T-HL)\vec v = H\vec e_1\), where \( \vec e_1 = (1,0,\ldots,0)^\mathsf{T} \), then
\[
1 + \sum_{i=1}^N v_i\ell^i = \frac{\det(T+HL)}{\det(T-HL)}.
\]
\end{prop}

\begin{proof}
Let \( V \) be the matrix whose first row is given by \( \big(\ell,\ell^2,\ldots,\ell^N\big) \) and whose remaining entries are all zero. Then
\begin{align*}
1 + \sum_{i=1}^N v_i\ell^i &= \vec e_1^\mathsf{T}\big(I + V (T-HL)^{-1} H\big)\vec e_1 \\
 &= \det\big(I + V (T-HL)^{-1} H\big) =\frac{\det(T+H(V-L))}{\det(T-HL)}.
\end{align*}
Thus, we need to show that \( \det(T+H(V-L))=\det(T+HL) \), i.e., that the matrix \( T+HL \) can be obtained by row and column operations from \( T+H(V-L) \). We shall perform the following operations:
\begin{itemize}
    \item for \( i \) from \( N \) to \( 2 \) (starting with the last row and ending with the second row) add \( \ell \) multiple of the \( i \)-th row to the \( (i-1)\)-st one;
    \item for \( j \) from \( N-1 \) to \( 1 \) (starting with the penultimate column and ending with the first one) subtract \( \ell \) multiple of the \( j \)-th column from the \( (j+1)\)-st one.
\end{itemize}
Set \( L_{i,j} := I + \ell E_{i,j} \), where \( E_{i,j} \) is the matrix with all zero entries except for \( (i,j) \)-th entry, which is \( 1 \). Since \(  I - \ell E_{i,j} = L_{i,j}^{-1} \), performing the suggested operations, is of course, equivalent to the multiplication on the left by \( L_{1,2} \cdots L_{N-1,N} \) and on the right by \( L_{N-1,N}^{-1} \cdots L_{1,2}^{-1} \). That is, we would like to show that
\[
L_{1,2} \cdots L_{N-1,N} (T+H(V-L)) L_{N-1,N}^{-1} \cdots L_{1,2}^{-1} = T+HL.
\]

It is straightforward to verify that
\[
\begin{cases}
L_{1,2} \cdots L_{N-1,N} T &= T L_{1,2} \cdots L_{N-1,N}, \\ 
L_{1,2} \cdots L_{N-1,N} H &= H L_{N,N-1}\cdots L_{2,1}.
\end{cases}
\]
For instance, when \( N=3 \), one can readily see that
\begin{multline*}
\begin{pmatrix} 1&\ell&0 \\ 0&1&0 \\ 0&0&1 \end{pmatrix} \begin{pmatrix} 1&0&0 \\ 0&1&\ell \\ 0&0&1 \end{pmatrix} T = \begin{pmatrix}
t_0 & t_1 + \ell t_0 & t_2 + \ell t_1 + \ell^2 t_0 \\
0 & t_0 & t_1 + \ell t_0 \\
0 & 0 & t_0
\end{pmatrix} \\ = T \begin{pmatrix} 1&\ell&0 \\ 0&1&0 \\ 0&0&1 \end{pmatrix} \begin{pmatrix} 1&0&0 \\ 0&1&\ell \\ 0&0&1 \end{pmatrix}.
\end{multline*}
and
\begin{multline*}
\begin{pmatrix} 1&\ell&0 \\ 0&1&0 \\ 0&0&1 \end{pmatrix}\begin{pmatrix} 1&0&0 \\ 0&1&\ell \\ 0&0&1 \end{pmatrix} H = \begin{pmatrix} h_1 + \ell h_2 + \ell^2 h_3 & h_2 + \ell h_3 & h_3 \\ h_2 + \ell h_3 & h_3 & 0 \\ h_3 & 0 & 0 \end{pmatrix} \\ = H \begin{pmatrix} 1&0&0 \\ 0&1&0 \\ 0&\ell&1 \end{pmatrix} \begin{pmatrix} 1&0&0 \\ \ell&1&0 \\ 0&0&1 \end{pmatrix}.
\end{multline*}
Hence, to prove the proposition, we need to show that
\[
L_{N,N-1}\cdots L_{2,1} (V-L) L_{N-1,N}^{-1} \cdots L_{1,2}^{-1} = L.
\]
Observe that
\[
L_{N,N-1}\cdots L_{2,1} V  = \big[\ell^{i+j-1}\big]_{i,j=1}^N.
\]
Therefore, the desired equality can be rewritten as
\[
\big[\ell^{i+j-1}\big]_{i,j=1}^N = L_{N,N-1}\cdots L_{2,1} L + L L_{1,2}\cdots L_{N-1,N}.
\]
The first summand on the right-hand side of the equality above represents successive additions of the \( (i-1) \)-st row multiplied by \( \ell \) to the \( i\)-th one starting with \( i=2 \). Thus, its \( (i,j) \)-th entry is given by
\begin{align*}
    \sum_{k=1}^i \ell^{i-k}\left(\frac\ell2\right)^{k+j-1}(-1)^{k-1}\binom{-j}{k-1} & = \frac{\ell^{i+j-1}}{2^j} \sum_{k=0}^{i-1} \left(-\frac12 \right)^k \binom{-j}{k} \\
    & = \left(\frac\ell2\right)^{i+j-1}\sum_{k=0}^{i-1} \binom{i+j-1}{k},
\end{align*}
where we used the equality of the first and middle terms in \eqref{bimonial identity} with \( x=y=\tfrac12\), \( m=i-1 \), and \( n=j \). On the other hand, the second summand represents successive additions of the \( (j-1) \)-st column multiplied by \( \ell \) to the \( j\)-th one starting with \( j=2 \). Hence, its \( (i,j) \)-th entry is given by
\begin{align*}
    \sum_{k=1}^j \ell^{j-k}\left(\frac\ell2\right)^{k+i-1}(-1)^{i-1}\binom{-k}{i-1} & = \frac{\ell^{i+j-1}}{2^i} \sum_{k=0}^{j-1} (-1)^{i-1} \left(\frac12\right)^k \binom{-k-1}{i-1} \\
    & = \left(\frac\ell2\right)^{i+j-1} \sum_{k=0}^{j-1}\binom{i+j-1}{k},
\end{align*}
where we used the equality of the first and last terms in \eqref{bimonial identity} with \( x=y=\tfrac12\), \( m=j-1 \), and \( n=i \). Since
\[
\sum_{k=0}^{i-1} \binom{i+j-1}{k} + \sum_{k=0}^{j-1}\binom{i+j-1}{i+j-1-k} = \sum_{k=0}^{i+j-1}\binom{i+j-1}{k} = 2^{i+j-1},
\]
the claim of the proposition follows.
\end{proof}

\end{document}
